\section{Analisis Kebutuhan}

Dari analisis masalah di atas, diturunkan kebutuhan rancangan yang harus dipenuhi artefak yang dibedakan menjadi kebutuhan fungsional dan kebutuhan non-fungsional.

Kebutuhan fungsional (KF):
\begin{enumerate}[label=\textbf{KF\arabic*.}, leftmargin=*]
    \item Menambang himpunan aturan ABAC berbentuk konjungsi kondisi atribut dari log akses, yang merekonstruksi keputusan izin/tolak.
    \item Menangkap keterkaitan berpasangan antar-atribut saat membangkitkan kandidat aturan.
    \item Memperbarui kebijakan secara inkremental ketika distribusi data berubah, termasuk mendeteksi kapan perubahan (\textit{drift}) tersebut terjadi.
\end{enumerate}

Kebutuhan non-fungsional (KNF):
\begin{enumerate}[label=\textbf{KNF\arabic*.}, leftmargin=*]
    \item[\textbf{KNF1}] \textbf{(memori)} --- Jejak memori bersifat $O(D)$ dan tidak tumbuh terhadap ukuran populasi, agar layak pada perangkat \textit{edge}.
    \item[\textbf{KNF2}] \textbf{(efisiensi pembaruan)} --- Pembaruan berkala hemat waktu; operasi mahal (pembelajaran struktur) hanya dijalankan bila diperlukan.
    \item[\textbf{KNF3}] \textbf{(kualitas)} --- Kualitas kebijakan non-inferior terhadap algoritma baseline sekelas.
\end{enumerate}

Keterkaitan antara persoalan, kebutuhan, dan keputusan rancangan yang menjawabnya dirangkum pada Tabel \ref{tab:pemetaan-masalah-kebutuhan}.
Dengan demikian, setiap komponen rancangan CoDA-EDA bukan merupakan pilihan sembarangan, melainkan jawaban langsung atas kebutuhan yang diturunkan dari analisis masalah.

\begin{table}[htbp]
\centering
\footnotesize
\setlength{\tabcolsep}{4pt}
\renewcommand{\arraystretch}{1.15}
\caption{Pemetaan Analisis Masalah --- Kebutuhan --- Keputusan Rancangan}
\label{tab:pemetaan-masalah-kebutuhan}
\begin{tabular}{@{} 
  >{\raggedright\arraybackslash}p{3.2cm} 
  >{\raggedright\arraybackslash}p{3.4cm} 
  >{\raggedright\arraybackslash}p{7.5cm} 
@{}}
\toprule
Persoalan & Kebutuhan & Keputusan rancangan \\
\midrule
1. Sumber daya \textit{edge} terbatas 
  & KNF1 (memori $O(D)$, tanpa populasi) 
  & \textit{Probability Vector} (PV) tanpa populasi eksplisit; pembaruan kompetitif $\pm 1/NP$ \\[2pt]

2. Dependensi vs.\ hemat memori 
  & KF2 + KNF1 
  & Pohon Chow-Liu orde-1 ($O(D)$), bukan jaringan dependensi penuh ($O(D^2)$) \\[2pt]

3. Dinamika \& biaya pembaruan 
  & KF3 + KNF2 
  & Pemisahan ritme: struktur dipicu-\textit{drift} (JSD $> \tau$) + parameter tiap siklus + \textit{warm-start} \\[2pt]

4. Kompromi kualitas proporsional 
  & KNF3 
  & Fungsi \textit{fitness} selaras-F1 (cakupan + presisi) + penambangan \textit{separate-and-conquer} \\
\bottomrule
\end{tabular}
\end{table}